\section{Toward Bethlehem}

\subsection{The House of Bread}\label{sec:bread}

Advent walks toward a town, and the prophet Micah named it:

\begin{quote}
And thou Bethlehem Ephrata, art a little one among the thousands of Juda, out
of thee shall he come forth unto me that is to be the ruler in Israel: and his
going forth is from the beginning, from the days of eternity.\par
{\smallskip\noindent\small ---Micah 5:2\par}
\end{quote}

A little one among the thousands, and from it one whose going forth is from
eternity. St Bernard, preaching on the Vigil, marvelled at it: ``It is not
in the royal city of Jerusalem, but in Bethlehem, which is the least of the
thousands of Juda. O Bethlehem! O little Bethlehem! once little, now magnified
by the Lord! He has magnified thee Who, though great, became little in thee.''

The name of the town means something. At Matins of Christmas, in the third nocturn,
the Church reads St Gregory the Great on St Luke's Gospel, and after a few lines
on the census he turns to it:

\begin{quote}
Then, the Lord is born in Bethlehem. Now the name Bethlehem signifieth ``the
House of Bread,'' and thus it is the birth-place of Him Who hath said, ``I am the
Living Bread, Which came down from heaven.'' (John vi. 51.) We see then that this
name of Bethlehem was prophetically given to the place where Christ was born,
because it was there that He was to appear in the flesh by Whom the souls of
the faithful are fed unto life eternal.\par
{\smallskip\noindent\small ---St Gregory the Great, Homily 8 on the Gospels\par}
\end{quote}

On the Vigil the Church took her Introit from Moses' promise of the manna. Our
Lord himself set the manna beside his own flesh, when the crowd at Capharnaum
reminded him that their fathers had eaten bread from heaven in the desert:

\begin{quote}
I am the bread of life. Your fathers did eat manna in the desert: and are dead.
This is the bread which cometh down from heaven: that if any man eat of it, he
may not die. I am the living bread which came down from heaven. If any man eat
of this bread, he shall live for ever: and the bread that I will give is my
flesh, for the life of the world.\par
{\smallskip\noindent\small ---John 6:48--52\par}
\end{quote}

\subsubsection*{Laid in a manger}

When he was born, his mother ``wrapped him up in swaddling clothes and laid him
in a manger''. On the first night of Advent the Church read of the ass that
knows his master's crib. Now the crib holds the Master himself, and he is laid
there as food. St Cyril of Alexandria writes:

\begin{quote}
He finds man in his corrupt affections become like the beasts that perish, and
therefore He is laid in the manger, in the place of food, that we changing the
life of beasts, might be brought to the knowledge that befits man, partaking
not of hay, but of the heavenly bread, the lifegiving body.\par
{\smallskip\noindent\small ---St Cyril of Alexandria, on St Luke\par}
\end{quote}

St Gregory the Great, in the same Christmas homily, takes the thought further,
from the hay in the manger to the hay that we are. Isaiah says, ``All flesh
is grass''; and Christ, made man, \latin{fenum nostrum vertit in frumentum}: he
turned our grass into wheat, he who said of himself that unless the grain of
wheat falls into the ground and dies, it remains alone. The Lord spoke of the
grain in the days before his Passion. Gregory hears it at the manger, where the
Child is laid to feed the faithful with the grain of his flesh. St Bede took the
thought from him: ``He Who is the bread of Angels is laid down in a manger, that
He might feast us, as it were the sacred animals, with the bread of His
flesh.''

\subsubsection*{From the manger to the altar}

St Ambrose, writing on the same verse, sets down a run of contrasts between
what Christ took and what he gave. Among them:
\latin{ille in praesaepibus, ut tu in altaribus}, ``he in the manger, that you
might be at the altars''. The manger of Bethlehem and the altar of the parish
church hold the same Lord, laid there for the same reason: that we may be
fed.

St Bernard turns this into a question about us. On the Vigil he tells his hearers
to make sure they are found in Bethlehem when the Lord is born there:

\begin{quote}
Remember, too, that it is in Bethlehem of Juda that Jesus is born; and be very
careful lest you fail to be found there, lest He fail to be received by you.
Bethlehem is the house of bread; Juda signifies confession or praise. If, then,
you replenish your soul with the food of the Divine Word, the Body of our Lord
Jesus Christ, and devoutly receive the Bread which came down from heaven, and
which giveth life to the world; if the vessel of your body is made strong and
able to hold the new wine by being refreshed and strengthened with His new and
glorified flesh; if, moreover, you live by faith, and have no need to weep
because you have forgotten to eat your bread, then, indeed, you are become a
Bethlehem fitted to receive our Lord.\par
{\smallskip\noindent\small ---St Bernard, First Sermon on the Vigil of the Nativity\par}
\end{quote}

St Bede, on the same Gospel, writes that ``down to the very end of time, the
Lord ceases not to be conceived at Nazareth, to be born at Bethlehem, whenever
any of His hearers taking of the flour of the word makes himself a house of
eternal bread.''

Make your Christmas Communion the centre of the feast, and prepare for it in
these last days as for the most important thing you will do at Christmas. Go to
confession in good time, in the Ember week if you can. In the
week before Christmas read the sixth chapter of St John, the discourse on the
bread of life, slowly, a few verses at a time. When you come to the crib on
Christmas night, look at the manger and remember what it is for.
